\documentclass[12pt,a4paper]{article}
\usepackage[utf8]{inputenc}
\usepackage[spanish]{babel}
\usepackage{amsmath, amssymb, amsthm}
\usepackage{geometry}
\geometry{top=2.5cm, bottom=2.5cm, left=2.5cm, right=2.5cm}
\usepackage{listings}
\usepackage{xcolor}

\lstset{
    language=Python,
    basicstyle=\ttfamily\small,
    keywordstyle=\color{blue},
    stringstyle=\color{red},
    commentstyle=\color{green!50!black},
    showstringspaces=false,
    frame=single,
    breaklines=true
}

\title{\textbf{Resumen Profundo Clase 4: El Arte de la Aproximación Numérica}}
\author{Ing. Omar Cáceres \and Modelado y Simulación}
\date{}

\begin{document}

\maketitle

\section{El Fundamento: ¿Por qué y Cómo aproximar?}
El desafío central del modelado numérico consiste en calcular un valor estimativo y sumamente preciso ($\hat{I}$) de una integral definida $I = \int_{a}^{b} f(x) dx$ cuando la función $f(x)$ carece parcial o totalmente de una primitiva analítica viable (ej., integrales de funciones gaussianas de probabilidad).

La deducción fundamental del error entre este valor aproximado y el real recurre al \textbf{Lema del Valor Medio de la Intergral}. Plantea que, en esencia, si $f$ y $g$ son continuas en un espectro y $g(x) \geq 0$, siempre existirá un punto intermedio libre y arbitrario $\xi \in [a, b]$ donde:
\begin{equation}
\int_{a}^{b} f(x)g(x) dx = f(\xi) \int_{a}^{b} g(x) dx
\end{equation}
Con base a esto, evaluamos poligonales y calculamos el desvío exacto.

\section{Regla del Rectángulo (Punto Medio)}
Plantea resolver subintervalos usando siempre el centro de los mismos para modelar una figura constante recta pura (polinomio de Grado 0).
\begin{itemize}
    \item \textbf{Aproximación Celda Unitaria $[-h/2, h/2]$:} $\int_{-h/2}^{h/2} f(x) dx \approx h \cdot f(0)$
    \item \textbf{Fórmula compuesta ($n$ particiones):} $I_r = h \sum_{k=0}^{n-1} f(x_{k+1/2})$
    \item \textbf{Error Global Teórico (Truncamiento):} $E_r = \frac{(b-a) h^2}{24} f''(\xi)$.
\end{itemize}

\section{Regla del Trapecio}
Abandona la constante pura y prefiere empalmar mediante una recta oblicua (secante de Grado 1) los pares de valores que engloban a cada subintervalo continuo evaluado.

\begin{itemize}
    \item \textbf{Fórmula compuesta:} 
    \begin{equation}
        I_t = \frac{h}{2} \left[ f(x_0) + 2\sum_{k=1}^{n-1} f(x_k) + f(x_n) \right]
    \end{equation}
    \item \textbf{Error Global Teórico:} $E_t = -\frac{(b-a)h^2}{12} f''(\xi)$.
\end{itemize}

\textbf{Detalle Técnico Clave:} Una intuición visual sugiere que empalmar con un techo sesgado en trapecio es inherentemente "mejor" que cortar a la mitad usando cajas llanas cuadradas (rectángulo). Contradictoriamente, el factor limitante $\frac{1}{24}$ en el Rectángulo y $\frac{1}{12}$ en el Trapecio evidencian teóricamente que el modelo Rectangular suele producir la mitad del error neto bajo un $h$ equitativo, si bien ambos comparten el orden global de convergencia $\mathcal{O}(h^2)$.

\section{Regla de Simpson 1/3 (La Malla Parabólica)}
Asciende al nivel Cuadrático de Newton-Cotes. Entraña envolver tríos sucesivos de puntos $(x_i, x_{i+1}, x_{i+2})$ con un contorno de molde de parábola de subgrado 2.
\begin{itemize}
    \item \textbf{Restricción Metodológica Ineludible:} El conteo puro de subintervalos $n$ debe ser rigurosamente un número par.
    \item \textbf{Fórmula compuesta generalizada:}
    \begin{equation}
        I_s = \frac{h}{3} \left[ f(x_0) + 4\sum_{i \text{ impar}} f(x_i) + 2\sum_{i \text{ par}} f(x_i) + f(x_n) \right]
    \end{equation}
    \item \textbf{Orden Crítico de Convergencia:} $\mathcal{O}(h^4)$. Un recorte de $h$ por dos contrae la varianza al valor de $\frac{1}{16}$ partes del error primario original. 
    \item \textbf{Error Global:} $E_{simp} = -\frac{(b-a)h^4}{180} f^{(4)}(\xi)$. \\
    \textit{Importante:} Al ser directamente dependiente aritméticamente de la 4ta Derivada general de la ecuación principal de base, este método modela y otorga exactitud pura y neta a las funciones primitivas cúbicas plenas convencionales ($f(x) \propto x^3$).
\end{itemize}

\section{Regla de Simpson 3/8 (Cúbica)}
Se expande interpolando no 3, sino 4 nodos con un polinomio cúbico puro (Grado 3).
\begin{itemize}
    \item \textbf{Restricción:} $n \pmod 3 \equiv 0$ ($n$ es un múltiplo de $3$).
    \item \textbf{Fórmula por Segmento:} $\approx \frac{3h}{8} \left[ f(a) + 3f(x_1) + 3f(x_2) + f(b) \right]$
    \item \textbf{Error Global Teórico:} Exactamente de orden $\mathcal{O}(h^4)$, cuantificado operativamente en $-\frac{(b-a)h^4}{80} f^{(4)}(\xi)$.
\end{itemize}

\section{Estimación Estructural de Ingeniería: Parámetro $h_{max}$}
No podemos operar asumiendo límites intangibles. Si imponemos o requerimos por contrato de modelo una \textit{Tolerancia Estructural Maximizada} tope equivalente a $\tau$ (p.ej., $\tau = 10^{-5}$):

Para Simpson 1/3, si asilamos y acotamos $\xi$ ubicando el peor de los casos máximos de nuestra respectiva cuota de derivada, definiéndolo como $M_4 = \max_{x \in [a,b]} |f^{(4)}(x)|$:

\begin{equation}
    |E| = \frac{b-a}{180} \cdot h^4 \cdot M_4 \leq \tau
\end{equation}
Despejamos el parámetro precalculable:
\begin{equation}
    h_{max} = \sqrt[4]{\frac{180 \cdot \tau}{(b-a) \cdot M_4}}
\end{equation}
Al configurar programáticamente una partición unánime $h \leq h_{max}$ nuestro script estará provisto de una garantía absoluta paramétrica asumiendo flotantes teóricos.

\section{Dilema Analítico: Truncamiento Matemático vs Limitante Físico de Redondeo}
La cota anterior presuponía operar calculadoras irreales teóricas matemáticas de precisión $\infty$.
Al traspasar los resultados al microprocesador:
\begin{itemize}
    \item \textbf{Error de Truncamiento}: Se acopla a las ecuaciones $\mathcal{O}(h^k)$ presentadas, mermando vertiginosamente al reducir dramáticamente el paso $h$.
    \item \textbf{Error de Redondeo Finito Flotante}: Es inversamente proporcional al tamaño de $h$. Interpolar con cortes infinitesimales minúsculos ($h \to 0$) infiere procesar brutalmente sumatorias en bucle por ciclos for. Cada una de ellas infiere e inyecta diminutas podredumbres de falta de memoria (coma flotante en bytes acotados). Hay un "Valle Óptimo" a partir del cual refinar $h$ degrada irreversiblemente a todo el cálculo puro final.
\end{itemize}

\section{Implementación en Script Puro de Python (NumPy Array Ops)}
\begin{lstlisting}
import numpy as np

# Configurar hiper-parametros
n = 120 # Subintervalos elegidos
h = (b - a) / n
x = np.linspace(a, b, n+1)
y = func(x)

# Regla del Trapecio
int_t = (h / 2) * (y[0] + 2 * np.sum(y[1:-1]) + y[-1])

# Regla de Simpson 1/3 (Garantizado un N = par localemente)
# Extrae el slice desde el index 1 y da saltos de a 2.
int_s = (h / 3) * (y[0] + y[-1] + 4 * np.sum(y[1:-1:2]) + 2 * np.sum(y[2:-2:2]))

# Regla de Simpson 3/8 (Garantizado un N = % 3)
int_s3 = (3 * h / 8) * (y[0] + y[-1] + 3 * np.sum(y[1:-1:3]) 
                        + 3 * np.sum(y[2:-1:3]) + 2 * np.sum(y[3:-2:3]))
\end{lstlisting}

\end{document}
